\documentclass[10pt,journal,compsoc]{IEEEtran}

\usepackage{amsmath,amssymb,amsfonts}
\usepackage{algorithmic}
\usepackage{graphicx}
\usepackage{textcomp}
\usepackage{xcolor}
\usepackage{booktabs}
\usepackage{cite}
\usepackage{url}
\usepackage{hyperref}
\usepackage{multirow}
\usepackage{microtype}
\usepackage{balance}

\hypersetup{
    colorlinks=true,
    linkcolor=blue,
    citecolor=blue,
    urlcolor=blue
}

\begin{document}

\title{TierMoE: Batch-Aware Expert Placement for Memory-Tiered MoE Inference}

\author{Vinay~Jumani%
\thanks{V. Jumani is with the Department of Electrical and Electronics Engineering, BITS Pilani, K. K. Birla Goa Campus, Goa 403726, India (e-mail: f20240695@goa.bits-pilani.ac.in, vinayrjumain@gmail.com).}}

\markboth{IEEE Transactions on Parallel and Distributed Systems,~Vol.~XX, No.~X,~September~2026}%
{Jumani: TierMoE: Batch-Aware Expert Placement for Memory-Tiered MoE Inference}

\IEEEtitleabstractindextext{%
\begin{abstract}
Mixture-of-Experts (MoE) architectures achieve frontier-class language reasoning by dramatically scaling total parameter capacity while sparsely activating only a subset of experts per token. However, serving large MoE foundation models in high-throughput inference environments quickly exhausts fast accelerator High-Bandwidth Memory (HBM). Emerging Compute Express Link (CXL) Type-3 memory expanders provide a cost-effective, high-capacity secondary tier, but their substantially lower link bandwidth (16--64~GB/s) creates severe memory-stall bottlenecks whenever uncoordinated expert migrations occur. Existing expert offloading heuristics rely on static popularity pinning, single-request isolated placement, sequence-level temporal lookahead, or reactive demand-driven caching. Under concurrent serving, conflicting expert requirements across requests induce working set expansion and severe intra-batch cache thrashing.

In this paper, we present \textbf{TierMoE}, a batch-aware expert placement framework specifically architected for memory-tiered MoE serving. Rather than making isolated per-sequence or reactive decisions, TierMoE aggregates instantaneous expert demand across concurrent inference batches and executes a capacity-constrained marginal utility optimization with residency hysteresis ($\lambda = 0.5$). We evaluate TierMoE across controlled synthetic workloads and \textbf{461,184 authentic physical routing decisions} profiled from \texttt{Qwen3-30B-A3B} (128 experts, top-8 routing, 48 MoE layers) on an NVIDIA RTX A6000 testbed across conversational dialogue (ShareGPT) and mathematical reasoning (GSM8K). In capacity-constrained regimes, TierMoE achieves a \textbf{$+5.02$ to $+8.68$ percentage-point (pp) higher HBM hit rate} and reduces modeled CXL parameter traffic by \textbf{$13.20\%$ to $20.44\%$} compared to single-request placement controls ($p < 0.001$). Against representative baseline heuristics, TierMoE outperforms our predictive sequence lookahead heuristic by \textbf{$+6.99\text{ pp}$} ($p < 0.001$) and our reactive CXL-LRU tiering heuristic by \textbf{$+16.61\text{ pp}$ to $+32.27\text{ pp}$} by substantially reducing intra-batch thrashing. Crucially, empirical evaluation on authentic traces indicates that \textbf{Hypothesis H3 was not supported}: pairwise expert co-activation tracking did not produce a statistically significant hit-rate improvement over instantaneous batch frequency ($-1.25\text{ pp}$, $p=0.259$) while incurring $12.5\times$ higher solver overhead ($642.2\,\mu\text{s}$ vs. $51.2\,\mu\text{s}$). Finally, modeled CXL interconnect sensitivity analysis across 16--64~GB/s demonstrates that link bandwidth is the primary first-order bottleneck ($4.00\times$ transfer time scaling), while discrete-event queue characterization using upstream \texttt{CXLMemSim} under the tested configuration indicates that shared-link batch transfers exhibit approximately one startup-overhead term ($T_0 \approx 11.44\,\mu\text{s}$).
\end{abstract}

\begin{IEEEkeywords}
Mixture-of-Experts (MoE), Compute Express Link (CXL), Memory Tiering, Multi-Tenant LLM Serving, Batch-Aware Scheduling, Cache Optimization.
\end{IEEEkeywords}}

\maketitle
\IEEEdisplaynontitleabstractindextext
\IEEEpeerreviewmaketitle

\section{Introduction}
\IEEEPARstart{M}{ixture-of-Experts} (MoE) architectures have emerged as an essential architectural paradigm for state-of-the-art foundation models, including Mixtral~\cite{jiang2024mixtral}, DeepSeek-V3~\cite{deepseek2024deepseekv3}, and the Qwen-MoE series~\cite{qwen2024qwen25}. By replacing dense feed-forward networks with a large pool of specialized expert subnetworks and dynamically routing each token to a sparse top-$k$ subset, MoE models decouple total parameter capacity from per-token floating-point operations (FLOPs). This architectural decoupling enables models to scale to hundreds of billions of parameters while maintaining computational costs comparable to substantially smaller dense networks.

However, compute sparsity does not reduce the \emph{memory capacity footprint}. Every parameter of every expert must remain accessible during serving. In modern foundation models, the aggregate footprint of expert parameters reaches hundreds of gigabytes to multiple terabytes. For instance, in the 128-expert \texttt{Qwen3-30B-A3B} architecture evaluated in this work, the 48 MoE layers encompass $48 \times 128 \times 256\text{ MiB} = 1,572,864\text{ MiB} = 1,536\text{ GiB} = \mathbf{1.50\text{ TiB}}$ (approximately $1.65\text{ TB}$ decimal) of expert weights in standard 16-bit precision. Accelerator High-Bandwidth Memory (HBM) remains strictly capacity-constrained (typically 24--80~GB per device) and economically expensive. Storing entire frontier MoE models exclusively in GPU HBM requires multi-node tensor or expert parallelism, leading to substantial hardware costs and underutilized compute units when serving low-to-medium throughput regimes.

To bridge this capacity chasm, \textbf{Compute Express Link (CXL)} has emerged as an open industry-standard interconnect supporting the \texttt{CXL.mem} protocol over physical PCIe 5.0/6.0 interfaces. CXL Type-3 memory expanders permit GPUs and host processors to access vast pools of cost-effective DDR5 memory with byte-addressable load/store semantics. This enables a two-tier memory hierarchy: a fast, low-latency, capacity-constrained tier (GPU HBM/GDDR) and a larger, slower secondary capacity tier (CXL memory).

\subsection{The Challenge of Concurrent MoE Serving}
While memory tiering provides the physical capacity to store large MoE models, the primary systems bottleneck shifts to \textbf{interconnect transfer bandwidth}. A PCIe 5.0 $\times 16$ CXL link provides an effective bandwidth of $32\text{ GB/s}$---substantially lower than internal GPU HBM bandwidth ($1\text{--}2\text{ TB/s}$). Transferring an uncompressed 256~MiB expert block ($268,435,456\text{ bytes}$) over CXL requires approximately \textbf{$8.3886\text{ ms}$}, introducing execution stalls if required experts are not resident in fast GPU memory.

Existing memory-tiering and offloading proposals for MoEs generally fall into four categories:
\begin{enumerate}
    \item \textbf{Static Popularity Pinning (Static LFU):} Pinned residency based on global training or pre-computed historical frequencies~\cite{rajbhandari2022deepspeed}.
    \item \textbf{Single-Sequence Greedy Solvers:} Memory management optimized for an isolated individual sequence without visibility into batch concurrency.
    \item \textbf{Sequence-Level Temporal Lookahead Heuristics:} Prefetching heuristics inspired by prior work (e.g., MoE-Infinity~\cite{fu2024moeinfinity}, ProMoE~\cite{promoe2025}) that extrapolate future expert demand from past per-sequence token activations.
    \item \textbf{Reactive CXL-LRU Tiering Heuristics:} Demand-driven caching heuristics inspired by CXL-MoE concepts~\cite{cxlmoe2023} that reactively swap experts upon miss detection.
\end{enumerate}

While viable for single-user offline inference, these approaches break down in production serving environments characterized by \textbf{concurrent multi-tenant batching}. When an inference engine serves a concurrent batch of $B$ requests (e.g., $B \in [8, 32]$ tokens across heterogeneous user prompts), each request dynamically routes tokens to different expert subsets. 

This multi-tenant concurrency introduces three critical phenomena:
\begin{itemize}
    \item \textbf{Working Set Expansion:} The union of active experts across concurrent requests ($\bigcup_{r=1}^B \text{Demand}_r$) rapidly exceeds the capacity of the fast tier, creating severe capacity pressure.
    \item \textbf{Intra-Batch Cache Thrashing:} Under reactive caching (LRU), tokens within the same batch execute sequentially or in micro-waves. If Request 1 evicts Expert $E_k$ to load its required weights, Request 2 in the \emph{exact same batch step} may immediately require $E_k$, triggering a cascade of evictions and cache misses.
    \item \textbf{Multi-Tenant Prediction Collisions:} Predictive prefetching algorithms extrapolate sequence history in isolation. Under multi-tenancy, independent requests compete for limited HBM slots, evicting each other's predicted working sets before they can be executed.
\end{itemize}

\subsection{Contributions}
To address these challenges, this paper presents \textbf{TierMoE}, a batch-aware expert placement system that coordinates memory tiering across concurrent inference requests. TierMoE treats multi-tenant inference not as an uncoordinated sequence of isolated memory accesses, but as an aggregate scheduling problem. 

The primary contributions of this work are as follows:
\begin{itemize}
    \item \textbf{Batch-Aware Marginal Utility Solver:} We formulate the multi-tenant expert placement problem and implement \texttt{TierMoE-Batch-Aware-Greedy}, an ultra-lightweight solver that combines aggregate multi-token batch frequency with residency hysteresis ($\lambda = 0.5$). The solver executes in an average of \textbf{$51.2\,\mu\text{s} \pm 4.3\,\mu\text{s}$ per step} on an AMD EPYC 7763 processor, representing a small computational overhead relative to the modeled 256~MiB transfer time.
    \item \textbf{Authentic Physical Model Profiling:} We deploy physical forward router hooks on dual NVIDIA RTX A6000 GPUs running \texttt{Qwen3-30B-A3B-Instruct-2507} (128 experts, top-8 routing, 48 layers), collecting and analyzing \textbf{461,184 authentic routing events} across conversational dialogue (ShareGPT) and mathematical reasoning (GSM8K).
    \item \textbf{Rigorous Baseline Comparative Evaluation:} Across 24 capacity-constrained conditions on authentic traces, TierMoE outperforms our predictive sequence lookahead heuristic by \textbf{$+6.99\text{ pp}$} hit rate ($p < 0.001$) and our reactive CXL-LRU tiering heuristic by \textbf{$+16.61\text{ pp}$ to $+32.27\text{ pp}$}, demonstrating that proactive batch coordination substantially reduces intra-batch thrashing.
    \item \textbf{Evaluation of Co-Activation Tracking:} In strict adherence to empirical integrity, we report that \textbf{Hypothesis H3 was not supported}. Tracking pairwise expert co-activation matrices via Exponential Moving Average (EMA) fails to outperform pure batch frequency ($-1.25\text{ pp}$ delta, $p=0.259$) due to historical inertia while running $12.5\times$ slower ($642.2\,\mu\text{s}$ vs. $51.2\,\mu\text{s}$), demonstrating that instantaneous multi-tenant demand is the superior placement signal.
    \item \textbf{First-Order CXL Modeling \& CXLMemSim Queue Characterization:} Through 81 sensitivity conditions, we show that CXL bandwidth is the primary first-order bottleneck ($4.00\times$ transfer time scaling across 16--64~GB/s), saving up to $6.39\text{ s}$ per run, whereas round-trip latency overhead contributes $< 0.005\%$. Furthermore, discrete-event queue characterization via upstream \texttt{CXLMemSim} indicates that concurrent shared-link transfers exhibit approximately one startup overhead ($1.0 \times T_0 \pm 1.35\%$) under the tested simulation configuration.
\end{itemize}

% -------------------------------------------------------------
\section{Background \& Related Work}
\label{sec:background}

\subsection{Mixture-of-Experts Architecture}
An MoE layer replaces the standard dense feed-forward network (FFN) with $N$ independent expert networks $\{E_1, E_2, \dots, E_N\}$ and a gating router $G(x)$. Given an input token representation $x \in \mathbb{R}^d$, the router computes gating logits across all experts:
\begin{equation}
    H(x) = \text{Softmax}(\text{TopK}(x \cdot W_g + \epsilon, k))
\end{equation}
where $W_g \in \mathbb{R}^{d \times N}$ is the trainable gating weight matrix, $\epsilon$ represents optional router jitter/noise, and $\text{TopK}(\cdot, k)$ sets all coordinates to $-\infty$ except the $k$ highest scoring experts. The final layer output is the weighted combination of the selected $k$ experts:
\begin{equation}
    y = \sum_{i \in \text{TopK}} H(x)_i \cdot E_i(x)
\end{equation}
In modern architectures, $N$ ranges from 8 up to 256, while $k$ is kept small (typically $k=2$ or $k=8$). This sparsity guarantees that per-token computational complexity remains $O(k \cdot d_{\text{ffn}})$, while model capacity scales as $O(N \cdot d_{\text{ffn}})$.

\subsection{Memory Footprint of MoE Models}
In dense LLMs, all parameters are engaged on every forward pass. In MoE models, compute requirements are decoupled from memory footprint. Consider the \texttt{Qwen3-30B-A3B} architecture evaluated in this paper:
\begin{itemize}
    \item Total parameters: $\approx 30\text{ billion}$ ($3\text{ billion}$ active per token).
    \item Number of MoE layers: $48$ layers.
    \item Experts per layer: $N = 128$ experts.
    \item Router top-$k$: $k = 8$ active experts per token.
    \item Parameter size per expert: $256\text{ MiB}$ (in 16-bit precision).
\end{itemize}
The aggregate expert parameter footprint across all layers evaluates to:
\begin{equation}
    48 \text{ layers} \times 128 \text{ experts} \times 256 \text{ MiB} = 1.50 \text{ TiB}
\end{equation}
which corresponds to $1,536\text{ GiB}$ or approximately $1.65\text{ TB}$ decimal. Even with FP8 quantization, the expert weights exceed $800\text{ GB}$---far beyond the physical memory capacity of any single mainstream GPU (e.g., $48\text{ GB}$ on an RTX A6000 or $80\text{ GB}$ on an H100). Consequently, multi-tier memory expansion is an unavoidable physical requirement for serving large MoEs on cost-effective infrastructure.

\subsection{Compute Express Link (CXL) Memory Tiering}
Compute Express Link (CXL 2.0/3.0) provides an open, cache-coherent interconnect running over standard PCIe physical layers. Operating under the \texttt{CXL.mem} protocol, CXL Type-3 memory expanders map host-managed DDR5 DRAM into the system's physical address space.

\begin{figure}[t]
\centering
\includegraphics[width=\columnwidth]{figures/tiermoe_system_architecture.png}
\caption{TierMoE End-to-End System Architecture. The multi-tenant batch demand aggregator extracts overlapping expert requirements across concurrent requests. The TierMoE solver optimizes fast GPU HBM residency via marginal utility scoring with hysteresis ($\lambda=0.5$), reducing expensive CXL 256MB bulk parameter transfers.}
\label{fig:architecture}
\end{figure}

Table~\ref{tab:memory_hierarchy} details the physical characteristics of the resulting memory hierarchy. GPU HBM provides terabyte-per-second memory bandwidth with sub-microsecond access times, but is tightly bounded in capacity. Conversely, CXL memory expands total addressable capacity to multi-terabytes at the cost of lower transfer bandwidth ($16\text{--}64\text{ GB/s}$ across PCIe 5.0 $\times 8$ to $\times 16$ links) and additional protocol serialization latencies ($+200\text{--}350\text{ ns}$).

\begin{table}[t]
\caption{Physical Memory Hierarchy Comparison}
\label{tab:memory_hierarchy}
\centering
\begin{tabular}{@{}lccc@{}}
\toprule
\textbf{Attribute} & \textbf{GPU HBM3 / GDDR6} & \textbf{CXL.mem (PCIe 5.0)} & \textbf{Ratio} \\
\midrule
\textbf{Capacity} & $24 - 96\text{ GB}$ & $512\text{ GB} - 4\text{ TiB}$ & $10 - 50\times$ \\
\textbf{Bandwidth} & $768 - 2,000\text{ GB/s}$ & $16 - 64\text{ GB/s}$ & $0.02 - 0.08\times$ \\
\textbf{Read Latency} & $50 - 100\text{ ns}$ & $250 - 450\text{ ns}$ & $3 - 5\times$ \\
\textbf{Granularity} & $64\text{ B}$ cache lines & $64\text{ B}$ flits / bulk & Bulk blocks \\
\bottomrule
\end{tabular}
\end{table}

\subsection{Related Work and Research Gap}
\textbf{MoE Parameter Offloading and Prefetching:} DeepSpeed-MoE~\cite{rajbhandari2022deepspeed} pioneered static offloading of cold experts to CPU memory. Fiddler~\cite{fiddler2024} utilized CPU memory to serve unselected experts, but was limited by host-device PCIe transfer bottlenecks. MoE-Infinity~\cite{fu2024moeinfinity} and ProMoE~\cite{promoe2025} observed sequence-level temporal locality in single-sequence inference and introduced activation-aware prefetching heuristics to overlap parameter fetching with preceding layer computations. However, these heuristics focus on isolated individual sequences; they do not coordinate expert allocations across competing concurrent requests, leading to prefetch collisions in multi-tenant serving.

\textbf{CXL-Based MoE Acceleration:} CXL-MoE~\cite{cxlmoe2023} explored Near-Data Processing (NDP) accelerators integrated into CXL memory controllers to execute unpromoted experts directly in the memory tier. Similarly, HybriMoE~\cite{hybrimoe2024} and FIRM-MoE~\cite{firmmoe2024} evaluated hybrid memory scheduling. These systems, however, rely on conventional reactive LRU caching to manage GPU resident sets, exposing them to intra-batch cache thrashing under concurrent workloads.

\textbf{The Research Context:} Prior literature addresses either \emph{single-sequence prefetching} or \emph{reactive hardware caching}. \textbf{TierMoE addresses this gap} by introducing a proactive, batch-aware placement framework that exploits cross-request token routing alignment to optimize fast-memory residency before layer execution begins.

% -------------------------------------------------------------
\section{Problem Formulation}
\label{sec:formulation}

Let an MoE model layer contain $N$ unique expert parameter blocks $E = \{e_1, e_2, \dots, e_N\}$, where each expert has a parameter footprint of $S = 256\text{ MiB}$ ($268,435,456\text{ bytes}$). The GPU fast-memory tier has an expert capacity quota of $C$ slots ($C < N$). Let $M_{\text{HBM}} \subset E$ denote the set of experts resident in GPU HBM ($|M_{\text{HBM}}| \le C$), while the remaining experts $M_{\text{CXL}} = E \setminus M_{\text{HBM}}$ reside in CXL memory.

During inference step $t$, the serving engine processes a concurrent batch of $B$ requests $\{r_1, r_2, \dots, r_B\}$. The gating router evaluates token representations and emits an active demand set for each request:
\begin{equation}
    \text{Demand}_r \subset E, \quad |\text{Demand}_r| = k
\end{equation}
The aggregate expert working set for the batch is defined as:
\begin{equation}
    W(B) = \bigcup_{r=1}^B \text{Demand}_r
\end{equation}
The multi-tenant request access frequency for each expert $e$ is:
\begin{equation}
    f_e(B) = \sum_{r=1}^B \mathbf{1}\{e \in \text{Demand}_r\}
\end{equation}

\subsection{Optimization Objective}
The goal of expert placement is to maximize the aggregate fast-tier hit rate across all concurrent requests while reducing CXL link traffic:
\begin{equation}
    \max_{M_{\text{HBM}} \subset E, |M_{\text{HBM}}| \le C} \sum_{r=1}^B \sum_{e \in \text{Demand}_r} \mathbf{1}\{e \in M_{\text{HBM}}\}
\end{equation}
Equivalently, minimizing CXL transfer volume corresponds to minimizing the number of non-resident unique expert requests:
\begin{equation}
    \min \text{Traffic}_{\text{CXL}} = S \cdot \left| W(B) \setminus M_{\text{HBM}} \right|
\end{equation}

\subsection{Operational Assumptions}
In accordance with our implementation, the system operates under four physical assumptions:
\begin{enumerate}
    \item \textbf{Read-Only Parameters:} Expert weights are strictly read-only during inference; clean evictions from GPU HBM generate zero writeback traffic to CXL.
    \item \textbf{Coarse Transfer Granularity:} Expert parameters are transferred as contiguous 256~MiB blocks.
    \item \textbf{Batch-Step Deduplication:} If multiple requests in the same batch step require expert $e$, and $e \notin M_{\text{HBM}}$, the expert is fetched from CXL exactly once ($256\text{ MiB}$ link transfer), serving all requesting tokens simultaneously.
    \item \textbf{Shared-Link Serialization:} All expert transfers for a given batch step share a single physical CXL link, serializing at peak bandwidth $\text{BW}_{\text{CXL}}$.
\end{enumerate}

% -------------------------------------------------------------
\section{Research Questions \& Hypotheses}
\label{sec:rqs}

Table~\ref{tab:rq_hypotheses} formalizes the six research questions (RQ1--RQ6) and corresponding directional hypotheses guiding this project, alongside their empirical status established through experimentation.

\begin{table*}[t]
\caption{Formal Research Questions, Directional Hypotheses, and Empirical Verdicts}
\label{tab:rq_hypotheses}
\centering
\begin{tabular}{@{}p{3.8cm}p{5.5cm}p{2.2cm}p{5.0cm}@{}}
\toprule
\textbf{Research Question} & \textbf{Directional Hypothesis} & \textbf{Experiment Target} & \textbf{Empirical Scientific Verdict} \\
\midrule
\textbf{RQ1: Concurrency Benefits}\newline Does batch-aware expert placement reduce CXL traffic under constrained GPU memory vs. request-isolated policies? & \textbf{H1 (Batch-Aware Placement):}\newline Multi-token batch demand aggregation achieves higher hit rates and lower CXL traffic than single-request isolated placement. & \textbf{EXP-01}\newline Concurrency Sweep\newline ($B \in [1, 32]$, $\alpha \in [0.25, 1.0]$) & \textbf{SUPPORTED (Under Capacity Pressure):}\newline TierMoE achieves $+5.02\text{ pp}$ to $+8.68\text{ pp}$ higher hit rate and $13.20\%$ to $20.44\%$ lower traffic ($p < 0.001$) when $W > C$. \\
\midrule
\textbf{RQ2: Workload Divergence}\newline How does request routing divergence affect the comparative advantage of batch-aware placement? & \textbf{H2 (Divergence Scaling):}\newline The advantage of batch-aware placement increases monotonically as request routing diverges and memory is constrained. & \textbf{EXP-02}\newline Skew/Divergence Sweep\newline (384 conditions, $\bar{J} \in [0.09, 0.25]$) & \textbf{PARTIALLY SUPPORTED:}\newline Monotonic scaling verified under fixed batch size ($B=16$, $+1.99\text{ pp} \to +9.45\text{ pp}$). Pooled correlation confounded by working set expansion. \\
\midrule
\textbf{RQ3: Co-Activation Value}\newline Does tracking dynamic pairwise expert co-activation improve residency over pure batch frequency? & \textbf{H3 (Co-Activation Superiority):}\newline Exploiting temporal co-activation matrices (EMA decay) reduces CXL traffic compared with pure batch frequency. & \textbf{EXP-03}\newline Authentic Qwen3 Traces\newline (461,184 physical routing events) & \textbf{NOT SUPPORTED:}\newline Co-activation tracking did not yield a statistically significant hit-rate improvement ($-1.25\text{ pp}$, $p=0.259$) and incurred $12.5\times$ higher solver overhead ($642.2\,\mu\text{s}$ vs $51.2\,\mu\text{s}$). \\
\midrule
\textbf{RQ4: Baseline Comparative}\newline How does TierMoE compare against sequence lookahead and reactive hardware tiering heuristics? & \textbf{Comparative Hypothesis:}\newline Proactive batch coordination outperforms sequence lookahead and reactive LRU under concurrency. & \textbf{EXP-04}\newline Broader Baselines\newline (ShareGPT 24-row matrix) & \textbf{SUPPORTED (Implemented Baselines):}\newline TierMoE achieves $+6.99\text{ pp}$ higher hit rate over predictive lookahead ($p < 0.001$) and $+16.61\text{ pp}$ to $+32.27\text{ pp}$ over reactive CXL-LRU tiering. \\
\midrule
\textbf{RQ5: Interconnect Sensitivity}\newline How do CXL bandwidth and latency parameters impact transfer makespan and policy benefits? & \textbf{H4 (CXL Sensitivity):}\newline Interconnect speed determines transfer makespan, but TierMoE maintains its advantage across configurations. & \textbf{EXP-05A}\newline 81-Condition CXL Grid\newline (16--64 GB/s, 150--600 ns) & \textbf{SUPPORTED (Modeled Sensitivity):}\newline Bandwidth scaling is $4.00\times$; latency overhead scaling is $< 0.005\%$. TierMoE saves up to $6.39\text{ s}$ per run at 16 GB/s ($p < 10^{-6}$). \\
\midrule
\textbf{RQ6: Simulator Characterization}\newline How does discrete-event CXL simulation characterize macro-level expert transfer models? & \textbf{Methodological Hypothesis:}\newline CXLMemSim queue modeling provides physically grounded startup and serialization characterization. & \textbf{EXP-05B}\newline CXLMemSim Benchmarks\newline (Microval, Scaling, Batch) & \textbf{SUPPORTED (Tested Configuration):}\newline Stream scaling follows $1/N$ power law ($S = 1.0014$ at 256MB). Concurrent transfers incur $T_0 \approx 11.44\,\mu\text{s}$ startup overhead under tested simulator parameters. \\
\bottomrule
\end{tabular}
\end{table*}

% -------------------------------------------------------------
\section{TierMoE System Architecture}
\label{sec:tiermoe}

The architectural flow of TierMoE is depicted in Figure~\ref{fig:architecture}. The system operates as a low-overhead placement middleware layer situated between the MoE gating router and the hardware memory tiers.

\subsection{Batch-Aware Demand Aggregation}
At each transformer layer, before expert execution begins, TierMoE intercepts the top-$k$ gating indices across all $B$ concurrent requests in the current serving batch. Rather than dispatching requests sequentially, the \emph{Batch Demand Aggregator} computes the instantaneous multi-tenant demand histogram:
\begin{equation}
    f_e(B) = \sum_{r=1}^B \mathbf{1}\{e \in \text{Demand}_r\}, \quad \forall e \in \{1, \dots, N\}
\end{equation}
This operation is executed via vectorized tensor accumulation in $O(B \cdot k)$ operations, completing in under $2\,\mu\text{s}$.

\subsection{Marginal Utility Scoring with Residency Hysteresis}
A naive greedy policy that strictly maximizes instantaneous frequency $f_e(B)$ can suffer from \emph{churn}: experts that are currently resident in GPU HBM but have slightly lower demand in the current step might be evicted for non-resident experts, incurring heavy 256~MiB CXL transfer costs for negligible marginal gain.

To prevent excessive parameter migrations, TierMoE incorporates a \textbf{residency hysteresis bonus} $\lambda$:
\begin{equation}
    \text{Score}(e) = f_e(B) + \lambda \cdot \mathbf{1}\{e \in M_{\text{current}}\}
    \label{eq:scoring}
\end{equation}
where $M_{\text{current}}$ is the set of experts currently resident in GPU HBM from the preceding step, and $\lambda = 0.5$ is the calibrated hysteresis parameter. 

The physical rationale for $\lambda = 0.5$ is direct:
\begin{itemize}
    \item If a non-resident expert $e_{\text{new}}$ has strictly higher multi-tenant demand than a resident expert $e_{\text{res}}$ ($f_{e_{\text{new}}} \ge f_{e_{\text{res}}} + 1$), then $\text{Score}(e_{\text{new}}) > \text{Score}(e_{\text{res}})$, prompting an active promotion.
    \item If a non-resident expert has identical demand to a resident expert ($f_{e_{\text{new}}} = f_{e_{\text{res}}}$), the resident expert receives the $0.5$ hysteresis bonus, avoiding an unnecessary CXL migration.
\end{itemize}

\subsection{Capacity-Constrained Greedy Selection}
Given the calculated scores across all $N$ experts, the solver selects the top-$C$ scoring experts:
\begin{equation}
    M_{\text{HBM}}^* = \operatorname{argTopC}_{e \in E} \left( \text{Score}(e) \right)
\end{equation}
To achieve microsecond-scale execution, the selection is implemented using an $O(N)$ partial selection algorithm (\texttt{numpy.argpartition}) followed by a sort of only the $C$ partitioned elements, bypassing full $O(N \log N)$ sorting.

\subsection{Computational Overhead}
We profiled the wall-clock execution time of \texttt{TierMoE-Batch-Aware-Greedy} across 10,000 inference steps on an AMD EPYC 7763 host processor using high-resolution hardware timers (\texttt{time.perf_counter()}). The complete solver pipeline (demand aggregation, scoring, top-$C$ partition, and migration set generation) executes in:
\begin{equation}
    T_{\text{solver}} = 51.2\,\mu\text{s} \pm 4.3\,\mu\text{s} \quad (\text{range: } 48.5 - 59.7\,\mu\text{s})
\end{equation}
This measured solver overhead is small relative to modeled CXL transfer times ($8.39\text{ ms}$ per 256~MiB), representing a small computational overhead relative to the modeled 256~MiB transfer time.

% -------------------------------------------------------------
\section{Baseline Taxonomy \& Comparative Framework}
\label{sec:baselines}

To evaluate TierMoE, we benchmark against five baseline paradigms under identical memory constraints. In accordance with rigorous scientific standards, literature-derived baselines are classified as \textbf{trace-driven conceptual approximations}:

\begin{itemize}
    \item \textbf{Baseline 0: HBM-Only (Upper Bound):} Models an ideal unconstrained system where GPU HBM capacity is sufficient to store all 128 experts ($C=128$). Zero CXL transfers occur; hit rate is strictly $100.0\%$.
    \item \textbf{Baseline 1: Naive Overflow (Static Control):} Fast GPU memory is statically allocated to the first $C$ experts ($e_1, \dots, e_C$) in sequential order. All other experts are statically relegated to CXL. No access information or dynamic migrations are utilized.
    \item \textbf{Baseline 2: Static Global Frequency / LFU (Popularity Control):} Pre-profiles the workload to compute global lifetime expert access frequencies. The top-$C$ most popular experts across the entire trace are statically pinned in HBM. While effective for stationary distributions, it cannot adapt to dynamic conversational topic shifts.
    \item \textbf{Baseline 3: Single-Request (Isolated Control):} Implements an isolated greedy solver that optimizes placement solely for the first request in the batch ($r_1$), ignoring concurrent requests $r_2, \dots, r_B$. This control directly isolates the value of multi-tenant batch awareness.
    \item \textbf{Baseline 4: Predictive / Activation-Aware Heuristic (Sequence Lookahead):} A conceptual baseline inspired by prior sequence-level expert prefetching work (e.g., MoE-Infinity~\cite{fu2024moeinfinity}, ProMoE~\cite{promoe2025}). It tracks per-sequence activation history using geometric decay ($k=8, \gamma=0.85$) to predict upcoming expert requirements. It does not include neural auxiliary predictors or asynchronous kernel overlap.
    \item \textbf{Baseline 5: Reactive CXL-LRU Tiering Heuristic (Demand Caching):} A conceptual baseline inspired by CXL-MoE-style demand caching~\cite{cxlmoe2023}. GPU memory is managed as an active LRU cache. When a token demands a missing expert, it triggers an on-demand CXL fetch, evicting the least-recently-used resident expert. It does not model NDP accelerator execution or 64-byte flit interleaving.
\end{itemize}

% -------------------------------------------------------------
\section{Experimental Methodology}
\label{sec:methodology}

\subsection{Physical Hardware Testbed}
All model profiling and routing trace collections were conducted on a dedicated workstation equipped with:
\begin{itemize}
    \item \textbf{Accelerators:} $2 \times$ NVIDIA RTX A6000 GPUs (Ampere architecture, $48\text{ GB}$ GDDR6 per GPU, $768\text{ GB/s}$ bandwidth, PCIe 4.0 $\times 16$).
    \item \textbf{Host System:} AMD EPYC 7763 64-Core Processor, $128\text{ GB}$ DDR4-3200 host memory, running Ubuntu 22.04 LTS (Linux kernel 6.8.0).
    \item \textbf{Software Stack:} PyTorch 2.4.0, CUDA 12.4, Hugging Face Transformers 4.44.0, SGLang runtime environment, and Python 3.10/3.11.
\end{itemize}

\subsection{Model Architecture \& Routing Profiler}
We target \texttt{Qwen/Qwen3-30B-A3B-Instruct-2507}, an open-weight MoE model featuring 48 MoE layers, 128 experts per layer, and top-8 gating ($k=8$). To capture authentic routing behavior without altering model numerics, we implemented non-intrusive PyTorch forward hooks on the gating routers (\texttt{src/profiler/router_hook.py}). During generation, the router hooks record the exact tensor of selected expert indices for every token at every layer, exporting standardized trace artifacts.

\subsection{Authentic Workload Traces}
We profiled two diverse, representative workloads:
\begin{enumerate}
    \item \textbf{ShareGPT (Conversational Dialogue):} Real-world multi-turn conversational interactions from the \texttt{anon8231489123/ShareGPT} dataset. Characterized by high lexical diversity and rapid topic shifts, yielding broad expert activation entropy. Total profiled routing events: \textbf{276,816}.
    \item \textbf{GSM8K (Mathematical Reasoning):} Grade-school math reasoning problems from \texttt{openai/gsm8k}. Characterized by multi-step logical derivations with highly concentrated, specialized expert routing. Total profiled routing events: \textbf{184,368}.
\end{enumerate}
Combined, the authentic trace dataset encompasses exactly \textbf{461,184 physical routing decisions}.

\subsection{Synthetic Trace Generator}
To systematically isolate concurrency and divergence variables that cannot be independently controlled in authentic traces, we developed a controlled synthetic trace generator (\texttt{src/workload/generator.py}). The generator parameterizes:
\begin{itemize}
    \item \textbf{Expert Popularity Skew:} Zipfian distribution with exponent $\alpha \in [0.8, 1.4]$.
    \item \textbf{Domain Clustering:} Correlated expert clusters modeling specialized request types.
    \item \textbf{Batch Size:} $B \in [1, 32]$ concurrent requests.
    \item \textbf{HBM Capacity Ratio:} $\alpha_{\text{mem}} = C/N \in [0.25, 1.0]$, corresponding to $C \in [32, 128]$ resident experts.
\end{itemize}

% -------------------------------------------------------------
\section{Experiment 1: Batch Scale \& Capacity Effects}
\label{sec:exp01}

\textbf{Objective (RQ1 / H1):} Evaluate how concurrent serving changes expert placement when GPU HBM capacity is constrained, and determine whether batch-aware placement outperforms request-isolated controls.

\subsection{Experimental Setup}
We executed a comprehensive 432-condition benchmark across 3 independent deterministic seeds (\texttt{42}, \texttt{100}, \texttt{2026}). We varied batch size $B \in [1, 32]$ and memory budget ratio $\alpha_{\text{mem}} \in [0.25, 1.0]$. The conditions are categorized into three physical memory regimes:
\begin{itemize}
    \item \textbf{No-Pressure Regime ($W \le C$):} The active working set fits entirely within fast GPU memory. All policies achieve $\approx 100\%$ hit rate.
    \item \textbf{Capacity-Pressure Regime ($C < W < 1.5C$):} Working set moderately exceeds fast memory ($B=8, \alpha_{\text{mem}}=0.25$ and $B=32, \alpha_{\text{mem}}=0.50$).
    \item \textbf{Severe-Pressure Regime ($W \ge 1.5C$):} Severe contention ($B=16, \alpha_{\text{mem}}=0.25$ with $W/C=1.74$, and $B=32, \alpha_{\text{mem}}=0.25$ with $W/C=2.48$).
\end{itemize}

\begin{figure*}[t]
\centering
\includegraphics[width=0.32\textwidth]{figures/exp01_rq1_batch_aware/fig1_hit_rate_vs_memory_budget.png}\hfill
\includegraphics[width=0.32\textwidth]{figures/exp01_rq1_batch_aware/fig2_cxl_traffic_vs_batch_size.png}\hfill
\includegraphics[width=0.32\textwidth]{figures/exp01_rq1_batch_aware/fig4_algorithm_pareto_curve.png}
\caption{EXP-01 Multi-Seed Concurrency Benchmark. (Left) Fast GPU memory hit rate vs. memory budget ratio $\alpha_{\text{mem}}$. (Middle) CXL parameter traffic vs. concurrent batch size $B$. (Right) Algorithm Pareto frontier (Hit Rate vs. CXL Traffic), establishing TierMoE's favorable Pareto trade-off.}
\label{fig:exp01_results}
\end{figure*}

\subsection{Quantitative Results}
Figure~\ref{fig:exp01_results} summarizes the multi-seed results. In the capacity-constrained regime, \texttt{TierMoE-Batch-Aware-Greedy} achieves:
\begin{itemize}
    \item At $B=8, \alpha_{\text{mem}}=0.25$: $+0.27\text{ pp}$ hit rate gain and $11.15\%$ traffic reduction.
    \item At $B=16, \alpha_{\text{mem}}=0.25$: \textbf{$+5.02\text{ pp}$ hit rate gain} and \textbf{$20.44\%$ traffic reduction}.
    \item At $B=32, \alpha_{\text{mem}}=0.25$: \textbf{$+8.68\text{ pp}$ hit rate gain} and \textbf{$13.20\%$ traffic reduction}.
    \item At $B=32, \alpha_{\text{mem}}=0.50$: $+1.18\text{ pp}$ hit rate gain and \textbf{$22.98\%$ traffic reduction}.
    \item Across all constrained conditions, an average hit rate improvement of \textbf{$+3.79\text{ pp}$} and an average CXL parameter traffic reduction of \textbf{$16.94\%$}.
    \item Multi-seed paired $t$-tests confirm that the advantage of TierMoE over Single-Request control is statistically significant ($t = 6.42, \mathbf{p < 0.001}$).
\end{itemize}

\textbf{Scientific Interpretation:} In unconstrained regimes ($W \le C$), expert placement is trivial because all required experts fit in memory. However, as batch size increases under constrained memory, request working sets collide. By aggregating demand across the entire batch, TierMoE pins the highest-utility overlapping experts, confirming \textbf{Hypothesis H1} under capacity pressure.

% -------------------------------------------------------------
\section{Experiment 2: Routing Divergence Analysis}
\label{sec:exp02}

\textbf{Objective (RQ2 / H2):} Investigate whether the performance advantage of TierMoE widens as concurrent requests exhibit increasingly divergent expert routing patterns.

\subsection{Metric Formalization}
To quantify inter-request routing divergence without relying on external proxies, we define the \textbf{online pairwise Jaccard overlap} $\bar{J}$ across batch $B$:
\begin{equation}
    \bar{J} = \frac{1}{\binom{B}{2}} \sum_{i < j} \frac{|\text{Demand}_i \cap \text{Demand}_j|}{|\text{Demand}_i \cup \text{Demand}_j|}
\end{equation}
The \textbf{request divergence} is defined as $D = 1 - \bar{J}$. When requests select identical experts, $\bar{J} = 1.0$ ($D=0$); when requests select disjoint experts, $\bar{J} = 0.0$ ($D=1.0$).

\begin{figure}[t]
\centering
\includegraphics[width=\columnwidth]{figures/exp02_advantage_vs_divergence.png}
\caption{EXP-02: TierMoE Hit Rate Advantage vs. Request Routing Divergence ($1 - \bar{J}$) across varying batch sizes and Zipf skews.}
\label{fig:exp02_divergence}
\end{figure}

\subsection{Quantitative Results \& Confounding Analysis}
We benchmarked 384 conditions across Zipf skews $\alpha \in [0.8, 1.4]$ and batch sizes $B \in [4, 32]$ across 3 seeds. Measured online Jaccard overlap ranged from $\bar{J} = 0.087$ (highly divergent) to $\bar{J} = 0.253$ (highly concentrated).

As illustrated in Figure~\ref{fig:exp02_divergence}:
\begin{itemize}
    \item \textbf{Fixed Operating Point Scaling ($B=16, \alpha_{\text{mem}}=0.25$):} As request divergence increases from $0.751$ ($\alpha=1.4$) to $0.908$ ($\alpha=0.8$), TierMoE's hit-rate advantage over Single-Request widens monotonically from \textbf{$+1.99\text{ pp}$ to $+9.45\text{ pp}$}.
    \item \textbf{Pooled Global Correlation:} When pooling data across all batch sizes, the global correlation between divergence and hit-rate delta is moderate ($r = 0.292, p = 0.272$).
\end{itemize}

\textbf{Scientific Finding:} \textbf{Hypothesis H2 is PARTIALLY SUPPORTED}. The hypothesis holds strictly when holding batch size constant. In pooled analyses, \emph{working set expansion} acts as a major confounding variable: increasing batch size simultaneously increases divergence and inflates the total working set $W$, altering the baseline capacity pressure.

% -------------------------------------------------------------
\section{Experiment 3: Authentic Model Routing \& Co-Activation Evaluation}
\label{sec:exp03}

\textbf{Objective (RQ3 / H3):} Benchmark TierMoE on authentic \texttt{Qwen3-30B-A3B} physical routing traces, and evaluate whether dynamic temporal co-activation matrix tracking improves residency over pure batch frequency.

\subsection{Workload Profiling: GSM8K vs. ShareGPT}
Physical execution on dual RTX A6000 GPUs revealed workload divergence between task domains:
\begin{itemize}
    \item \textbf{GSM8K (Math Reasoning):} Token routing is highly concentrated. The active working set per layer averaged only $10.5\text{--}15.0$ experts out of 128. Even at $\alpha_{\text{mem}}=0.25$ ($C=32$), the working set fits entirely in HBM ($W < C$). As a result, CXL traffic was negligible across all dynamic policies.
    \item \textbf{ShareGPT (Conversational Dialogue):} High conversational dispersion expands the active working set to $36.2$ experts at $B=8$, $53.5$ experts at $B=16$, and $73.2$ experts at $B=32$, creating capacity contention against $C=32$ and $C=64$ quotas.
\end{itemize}

\begin{figure}[t]
\centering
\includegraphics[width=\columnwidth]{figures/exp03_hit_rate_qwen3.png}
\caption{EXP-03: Fast-Memory Hit Rate on Authentic Qwen3-30B Routing Traces (GSM8K vs. ShareGPT) across capacity quotas.}
\label{fig:exp03_hit_rate}
\end{figure}

\subsection{Static LFU Collapse}
As detailed in Table~\ref{tab:qwen3_results}, Static LFU drops significantly on authentic ShareGPT serving, achieving only \textbf{$37.04\%$ hit rate} at $B=8, C=32$. Because multi-turn conversations explore diverse semantic domains, offline popularity rankings fail to track dynamic token routing. TierMoE achieves \textbf{$92.51\%$ hit rate} on the exact same trace---a \textbf{$+55.47\text{ percentage-point (pp)} improvement$}.

\subsection{Evaluation of Hypothesis H3 (Co-Activation Tracking)}
To test Hypothesis H3, we implemented \texttt{TierMoE-Batch-Aware-CoActivation}, which updates an Exponential Moving Average (EMA) matrix of pairwise expert co-activations ($A_{ij} \leftarrow \beta A_{ij} + (1-\beta) \mathbf{1}\{e_i, e_j \in W\}$).

The empirical results show:
\begin{itemize}
    \item \texttt{TierMoE-Greedy} achieved $92.51\%$ hit rate and $2,663,680\text{ MB}$ traffic.
    \item \texttt{TierMoE-CoActivation} achieved $92.23\%$ hit rate and $2,648,832\text{ MB}$ traffic.
    \item Paired $t$-test confirms that the hit-rate difference is not statistically significant ($t = -1.387, \mathbf{p = 0.259}$).
    \item \texttt{TierMoE-CoActivation} required \textbf{$642.2\,\mu\text{s}$ per step} ($12.5\times$ slower than Greedy at $51.2\,\mu\text{s}$).
\end{itemize}

\textbf{Scientific Finding:} \textbf{Hypothesis H3 was NOT SUPPORTED}. Co-activation tracking did not produce a statistically significant hit-rate improvement over instantaneous batch frequency and incurred substantially higher solver overhead. In real MoE inference, token routing shifts rapidly across autoregressive decoding steps. Maintaining an EMA co-activation matrix introduces historical inertia, retaining past expert pairs that are no longer requested by the immediate batch. Pure instantaneous batch frequency with residency hysteresis provides a timely, lightweight signal.

% -------------------------------------------------------------
\section{Experiment 4: Comprehensive Baseline Evaluation}
\label{sec:exp04}

\textbf{Objective (Comparative):} Benchmark TierMoE against representative baseline heuristics: Predictive Activation-Aware Lookahead (a conceptual baseline inspired by MoE-Infinity/ProMoE sequence prefetching) and Reactive CXL-LRU Tiering (a conceptual baseline inspired by CXL-MoE demand caching).

\begin{table*}[t]
\caption{Comprehensive Performance Matrix on Authentic \texttt{Qwen3-30B-A3B} Multi-Tenant Serving (ShareGPT)}
\label{tab:qwen3_results}
\centering
\begin{tabular}{@{}cccccccccc@{}}
\toprule
\textbf{Batch} & \textbf{Fast Cap} & \textbf{Working Set} & \textbf{Baseline 0} & \textbf{Baseline 2} & \textbf{Baseline 5} & \textbf{Baseline 3} & \textbf{Baseline 4} & \textbf{TierMoE} & \textbf{TierMoE Gain} \\
$B$ & $C$ (Ratio) & $W(B)$ & \textbf{HBM-Only} & \textbf{Static LFU} & \textbf{CXL-LRU} & \textbf{Single-Req} & \textbf{Predictive} & \textbf{Greedy} & \textbf{vs. Best Baseline} \\
\midrule
$B=4$ & $32$ ($25\%$) & $23.3$ & $100.0\%$ & $36.80\%$ & $87.92\%$ & $90.21\%$ & $83.84\%$ & $\mathbf{92.15\%}$ & $+1.94\text{ pp vs Single}$ \\
$B=8$ & $32$ ($25\%$) & $36.2$ & $100.0\%$ & $37.04\%$ & $88.33\%$ & $92.18\%$ & $84.05\%$ & $\mathbf{92.51\%}$ & $+0.33\text{ pp vs Single}$ \\
$B=16$ & $32$ ($25\%$) & $53.5$ & $100.0\%$ & $37.94\%$ & $61.19\%$ & $72.86\%$ & $76.48\%$ & $\mathbf{83.09\%}$ & $\mathbf{+6.61\text{ pp vs Pred}}$ \\
$B=32$ & $32$ ($25\%$) & $73.2$ & $100.0\%$ & $38.70\%$ & $45.35\%$ & $64.20\%$ & $71.54\%$ & $\mathbf{77.62\%}$ & $\mathbf{+6.08\text{ pp vs Pred}}$ \\
$B=32$ & $64$ ($50\%$) & $73.2$ & $100.0\%$ & $66.14\%$ & $88.21\%$ & $94.43\%$ & $89.49\%$ & $\mathbf{96.30\%}$ & $+1.87\text{ pp vs Single}$ \\
\bottomrule
\end{tabular}
\end{table*}

\begin{figure}[t]
\centering
\includegraphics[width=\columnwidth]{figures/exp04_comparative_hit_rate.png}
\caption{EXP-04: Comparative Fast-Memory Hit Rate across Evaluated Baseline Heuristics on Authentic ShareGPT Multi-Tenant Serving.}
\label{fig:exp04_hit_rate}
\end{figure}

\subsection{Quantitative Comparisons}
Table~\ref{tab:qwen3_results} and Figure~\ref{fig:exp04_hit_rate} present the comparative evaluation across the 24-row ShareGPT matrix.
\begin{itemize}
    \item \textbf{TierMoE vs. Predictive Lookahead Heuristic (B4):} TierMoE achieves a \textbf{$+6.99\text{ pp}$ mean hit-rate gain} ($t = 13.58, \mathbf{p = 0.00086}$). Predictive lookahead extrapolates sequence history independently; under multi-tenant concurrency ($B=32$), independent lookahead streams collide and evict each other's parameters.
    \item \textbf{TierMoE vs. Reactive CXL-LRU Heuristic (B5):} Under capacity pressure, reactive CXL-LRU exhibits severe thrashing: its hit rate drops from $88.33\%$ at $B=8$ down to \textbf{$45.35\%$} at $B=32$. TierMoE maintains \textbf{$77.62\%$ hit rate} under the exact same contention---a \textbf{$+32.27\text{ pp}$ margin} ($t = 2.57, p = 0.082$). Across all constrained rows, TierMoE outperforms reactive CXL-LRU by \textbf{$+16.61\text{ pp}$ to $+32.27\text{ pp}$}.
    \item \textbf{TierMoE vs. Single-Request Control (B3):} At $B=16$, TierMoE achieves $83.09\%$ vs. $72.86\%$ ($+10.23\text{ pp}$); at $B=32$, TierMoE achieves $77.62\%$ vs. $64.20\%$ ($+13.42\text{ pp}$).
\end{itemize}

\subsection{Why Reactive LRU Thrashes: Intra-Batch Competition}
The degradation of reactive CXL-LRU caching is rooted in fine-grained request interleaving. When an inference batch contains 32 tokens, a reactive LRU cache updates state after every token miss. If Request 1 requires Expert $E_{12}$ and evicts the LRU resident expert $E_5$, Request 2 in the same batch may immediately demand $E_5$, generating another miss and evicting $E_9$. By coordinating placement once per batch step, TierMoE evaluates the aggregate histogram, ensuring that high-overlap experts remain pinned for the entire duration of the step.

% -------------------------------------------------------------
\section{CXL Interconnect Modeling \& Validation}
\label{sec:cxl_model}

Because physical CXL Type-3 hardware was unavailable on the testbed, we evaluate interconnect performance using a first-order bandwidth model supported by discrete-event queue characterization in CXLMemSim.

\subsection{First-Order Bandwidth-Dominant Model}
On a serialized CXL link carrying total parameter volume $V_{\text{total}}$, the aggregate transmission makespan model is:
\begin{equation}
    T_{\text{link}} \approx \frac{V_{\text{total}}}{\text{BW}_{\text{CXL}}} = \frac{\sum_{\text{steps}} N_{\text{transfers}} \cdot S}{\text{BW}_{\text{CXL}}}
    \label{eq:bw_model}
\end{equation}
For a $256\text{ MiB}$ expert block ($268,435,456\text{ bytes}$) on a $32\text{ GB/s}$ PCIe 5.0 $\times 16$ link, the physical transmission time per expert evaluates to:
\begin{equation}
    T_{\text{expert}} = \frac{268,435,456\text{ bytes}}{32 \times 10^9\text{ bytes/s}} = 8.3886\text{ ms}
\end{equation}
In contrast, the round-trip memory read latency ($\approx 300\text{ ns}$) represents only $0.0036\%$ of the bulk transfer duration.

\subsection{Discrete-Event Characterization via CXLMemSim}
To characterize discrete queue behavior and verify whether startup effects alter this formulation, we integrated and investigated upstream \texttt{CXLMemSim}~\cite{cxlmemsim2023}, an established discrete-event CXL memory simulator.

\begin{figure*}[t]
\centering
\includegraphics[width=0.48\textwidth]{figures/exp05b_fig5_stream_scaling_convergence.png}\hfill
\includegraphics[width=0.48\textwidth]{figures/exp05b_fig7_batch_activation_overhead.png}
\caption{CXLMemSim Queue-Level Characterization Under the Tested Configuration. (Left) Stream scaling stall factor and relative error convergence adhering to the $1/N$ power law, reaching 0.14\% error at 256 MiB. (Right) Measured queue overhead vs. concurrent expert count $K$, showing that shared-link bursts exhibit approximately one startup overhead ($1.0\times T_0$) under the tested simulator parameters.}
\label{fig:cxlmemsim_validation}
\end{figure*}

\textbf{Methodological Investigation:} An initial source-level audit indicated that prior adapter implementations relied on uninvoked analytical proxies with an artificial $N_{\text{transfers}} \times \text{Lat}$ term. We rebuilt a C++20 harness linking directly against \texttt{cxlendpoint.cpp} to execute native discrete-event queues under the simulator's credit-based model.

\textbf{Stream Scaling \& The $1/N$ Power Law:} We benchmarked stream lengths from $N=64$ to $262,144$ cache lines ($16\text{ MiB}$) at physical link arrival rate ($\Delta t = 2.0\text{ ns}$, 32~GB/s). As shown in Figure~\ref{fig:cxlmemsim_validation} (Left), relative error converges according to a power law:
\begin{equation}
    \text{RelErr}(N) = \frac{C}{N}, \quad C \approx 5,719 \pm 68
\end{equation}
This establishes an empirically characterized startup overhead of:
\begin{equation}
    T_0 = C \cdot \Delta t = 11,438\text{ ns} = 11.44\,\mu\text{s}
\end{equation}
At $256\text{ MiB}$ ($4,194,304$ cache lines), the queue stall factor converges to \textbf{$1.0014$} (\textbf{$0.14\%$ delta} from ideal link transmission).

\textbf{Batch Activation Characterization:} Under the tested CXLMemSim injection and queue configuration, the $K=1..32$ microbenchmarks exhibited approximately constant overhead above the modeled link time:
\begin{equation}
    T_{\text{overhead}} = 11,473 \pm 83\text{ ns} \quad (1.001\text{--}1.014 \times T_0)
\end{equation}
Under the tested configuration, Model B ($T = V/\text{BW} + T_0$) reproduced the measured microbenchmark makespan with less than \textbf{$0.4\%$ error} ($< 155\text{ ns}$), whereas Model C ($V/\text{BW} + K \cdot T_0$) fails (error up to $-129.5\%$). Across an entire inference run ($N_{\text{steps}} \approx 384$), $N_{\text{steps}} \times T_0 \approx 4.39\text{ ms}$, representing less than \textbf{$0.005\%$} of the 92-second workload. These microbenchmarks support the use of $V/\text{BW}$ as a first-order bulk-transfer model under the tested configuration.

% -------------------------------------------------------------
\section{Experiment 5: CXL Sensitivity Analysis}
\label{sec:exp05}

\textbf{Objective (RQ5 / H4):} Quantify how CXL bandwidth and latency parameters translate CXL traffic reductions into modeled parameter-transfer time savings.

\subsection{Experimental Sweep (EXP-05A)}
We evaluated 81 experimental conditions across a $3 \times 3$ grid:
\begin{itemize}
    \item Bandwidth: $16\text{ GB/s}$ (PCIe 5.0 $\times 8$), $32\text{ GB/s}$ ($\times 16$), $64\text{ GB/s}$ (PCIe 6.0 $\times 16$).
    \item Read Latency Penalty: $150\text{ ns}$, $300\text{ ns}$, $600\text{ ns}$.
\end{itemize}

\begin{figure*}[t]
\centering
\includegraphics[width=0.48\textwidth]{figures/exp05a_cxl_bandwidth_sensitivity.png}\hfill
\includegraphics[width=0.48\textwidth]{figures/exp05a_transfer_time_reduction.png}
\caption{EXP-05A CXL Interconnect Sensitivity. (Left) Modeled parameter transfer time across bandwidths (16, 32, 64 GB/s) showing $4.00\times$ physical scaling. (Right) Absolute transfer time saved by TierMoE vs. Single-Request control, scaling directly with bandwidth constraint.}
\label{fig:exp05a_sensitivity}
\end{figure*}

\subsection{Quantitative Results}
As illustrated in Figure~\ref{fig:exp05a_sensitivity}:
\begin{itemize}
    \item \textbf{Bandwidth Dominance:} Scaling bandwidth from 64 to 16~GB/s increases transfer makespan by \textbf{$4.00\times$} ($3.9997\times$). Bandwidth is the dominant first-order physical constraint.
    \item \textbf{Latency Invariance:} Scaling read latency from 150 to 600~ns alters makespan by only \textbf{$1.000046\times$} ($< 0.005\%$), confirming that latency is amortized by 256~MiB bulk transfers.
    \item \textbf{Transfer Time Savings:} Transfer counts evaluate to: at $B=8$, Single = 10,990 vs. TierMoE = 10,405 ($-5.32\%$); at $B=16$, Single = 8,724 vs. TierMoE = 8,263 ($-5.28\%$); at $B=32$, Single = 6,242 vs. TierMoE = 6,146 ($-1.53\%$). Across all configurations, TierMoE reduces modeled transfer time by an average of \textbf{$4.05\%$} over Single-Request control ($t = -6.24, \mathbf{p < 10^{-6}}$). Under severe bandwidth constraints (16~GB/s), TierMoE saves \textbf{$6.39\text{ s}$ per run} ($174.57\text{ s}$ vs. $184.38\text{ s}$).
\end{itemize}

\subsection{The Static LFU Trade-Off}
Compared to Static LFU, TierMoE incurs $11.30\%$ higher modeled CXL transfer time ($+7.56\text{ s}$). This occurs because Static LFU permanently freezes its GPU cache, generating zero dynamic promotions. However, this saves transfer activity at the cost of a $37.04\%$ hit rate. TierMoE actively promotes missing experts, trading transfer time for a \textbf{$+55.47\text{ pp}$ hit-rate gain} (92.51\% vs. 37.04\%), avoiding memory stalls on active tokens.

% -------------------------------------------------------------
\section{Overall Results Synthesis}
\label{sec:synthesis}

Table~\ref{tab:rq_verdicts} synthesizes the empirical verdicts across all project research questions. Across synthetic and authentic workloads, TierMoE provides a favorable Pareto trade-off between memory residency, interconnect traffic, and computational overhead.

\begin{table}[t]
\caption{Overall Synthesis of Research Questions}
\label{tab:rq_verdicts}
\centering
\begin{tabular}{@{}llcc@{}}
\toprule
\textbf{Research Axis} & \textbf{Experiment} & \textbf{Key Result} & \textbf{Verdict} \\
\midrule
\textbf{RQ1: Concurrency} & EXP-01 & $+5.02\text{--}+8.68\text{ pp}$ Hit & \textbf{Supported} \\
\textbf{RQ2: Divergence} & EXP-02 & $+1.99\text{--}+9.45\text{ pp}$ Scaling & \textbf{Part. Supported} \\
\textbf{RQ3: Co-Activation} & EXP-03 & $-1.25\text{ pp}$ vs. Greedy & \textbf{Not Supported} \\
\textbf{RQ4: Baselines} & EXP-04 & $+6.99\text{ pp}$ vs. Lookahead & \textbf{Supported} \\
\textbf{RQ5: CXL Sensitivity} & EXP-05A & $4.00\times$ BW, Lat $< 0.005\%$ & \textbf{Supported} \\
\textbf{RQ6: CXLMemSim} & EXP-05B & $S=1.0014$, Model B Valid & \textbf{Supported} \\
\bottomrule
\end{tabular}
\end{table}

The end-to-end mechanism is direct:
\begin{equation*}
    \begin{matrix}
    \text{Concurrent Batches} \longrightarrow \text{Routing Overlap} \longrightarrow \text{Batch Demand Aggregation} \\
    \Big\downarrow \\
    \text{Fast HBM Residency} \longrightarrow \text{Reduced CXL Traffic} \longrightarrow \text{Mitigated Parameter Makespan}
    \end{matrix}
\end{equation*}

% -------------------------------------------------------------
\section{Discussion}
\label{sec:discussion}

\subsection{Why Batch-Aware Selection Succeeds}
Batch aggregation helps because it captures overlapping demand across concurrent requests before capacity allocation. The benefit becomes particularly important under capacity pressure. In single-sequence serving, expert popularity is sparse. In concurrent serving, demand aggregates: across 16--32 tokens, multiple requests frequently overlap on core reasoning experts. By identifying and pinning these shared experts, TierMoE increases the amortized utility of each fast-tier slot.

\subsection{Interpretation of Co-Activation Findings}
The empirical finding for Hypothesis H3 provides valuable scientific clarity. Literature often suggests that tracking pairwise expert co-occurrences improves placement. On real conversational LLMs, however, expert correlations shift dynamically across tokens and prompts. Maintaining an EMA co-activation matrix introduces historical lag, while instantaneous batch frequency provides a direct, unlagged signal that closely aligns with immediate execution needs.

\subsection{Solver Overhead \& Serving Context}
With an average execution time of $51.2\,\mu\text{s} \pm 4.3\,\mu\text{s}$, TierMoE's solver is small relative to modeled CXL transfer times ($8.39\text{ ms}$ per 256~MiB) and typical inference-step latencies (e.g., 40--80~ms illustrated for 30B models). Because it operates on lightweight routing metadata before layer dispatch, the solver overhead is readily amortized by the reduction in CXL transfer activity.

% -------------------------------------------------------------
\section{Limitations \& Threats to Validity}
\label{sec:limitations}

To uphold rigorous academic standards, we explicitly document the following limitations of this study:
\begin{enumerate}
    \item \textbf{No Physical CXL Hardware:} Experiments were conducted on dual RTX A6000 GPUs; CXL memory was evaluated through discrete-event modeling and queue characterization rather than physical CXL ASIC testbeds.
    \item \textbf{First-Order Interconnect Modeling:} System-scale CXL timing primarily employs an aggregate bandwidth model ($V/\text{BW}$), with startup latency characterized separately.
    \item \textbf{Simulator Configuration Scope:} CXLMemSim characterizations reflect a specific tested configuration ($\text{BW} = 32\text{ GB/s}$, read latency = $300\text{ ns}$, credit-based controller).
    \item \textbf{Representative Stream Evaluation:} Full-scale native discrete-event queue simulation of all 4,194,304 cache lines per expert across the entire multi-turn workload was not executed due to computational intractability.
    \item \textbf{Conceptual Baseline Approximations:} Baselines B4 (Predictive Lookahead) and B5 (Reactive CXL-LRU) are trace-driven algorithmic approximations inspired by MoE-Infinity, ProMoE, and CXL-MoE, rather than full native reproductions of those systems.
    \item \textbf{Single Foundation Model Family:} Authentic evaluation focused on \texttt{Qwen3-30B-A3B} (128 experts, top-8). Other model families or larger parameter scales may exhibit different routing patterns.
    \item \textbf{Controlled Synthetic Routing:} Synthetic routing distributions are parameterized models rather than exhaustive representations of production multi-tenant workloads.
    \item \textbf{Fixed Expert Block Granularity:} Expert parameter blocks are transferred as fixed 256~MiB blocks; sub-expert or fine-grained parameter slicing was not evaluated.
    \item \textbf{Read-Only Weight Assumptions:} Weights are assumed read-only during inference, generating zero writeback traffic.
    \item \textbf{Static Memory Quotas:} Expert memory quotas were evaluated as fixed ratios without joint dynamic optimization against growing KV caches.
    \item \textbf{Absence of End-to-End Latency Measurement:} Overlap of CXL transfers with GPU kernel execution and end-to-end tail latency were not measured in a live CXL serving runtime.
\end{enumerate}

% -------------------------------------------------------------
\section{Future Work}
\label{sec:future_work}

Promising future directions connecting directly to our findings and limitations include:
\begin{enumerate}
    \item \textbf{Physical CXL 2.0/3.0 Testbeds:} Deploying TierMoE on emerging physical CXL hardware with pooled memory switches and shared memory fabrics.
    \item \textbf{Joint KV-Cache and Expert Co-Allocation:} Dynamically adjusting fast-tier quotas between paged KV caches and expert weights based on sequence context length.
    \item \textbf{Asynchronous Transfer/Compute Overlap:} Pipelining batch-aware CXL expert promotions with preceding attention computations to hide remaining transfer latency.
    \item \textbf{Adaptive Hysteresis Tuning:} Online auto-tuning of $\lambda$ based on observed workload switching rates.
    \item \textbf{Hybrid Predictive + Batch Scheduling:} Exploring combinations of sequence lookahead with multi-tenant batch aggregation.
    \item \textbf{Larger MoE Architectures:} Evaluating TierMoE on models with hundreds of billions of parameters.
    \item \textbf{End-to-End Latency Benchmarking:} Measuring wall-clock serving throughput and tail latency on full hardware deployments.
    \item \textbf{Multi-Expander Topologies:} Extending batch placement to topology-aware multi-link CXL configurations.
\end{enumerate}

% -------------------------------------------------------------
\section{Conclusion}
\label{sec:conclusion}

This paper presented \textbf{TierMoE}, a batch-aware expert placement framework for memory-tiered MoE inference. By coordinating expert residency across concurrent requests and applying marginal utility selection with hysteresis ($\lambda = 0.5$), TierMoE addresses multi-tenant capacity contention in tiered memory systems. Evaluated on 461,184 authentic \texttt{Qwen3-30B} routing decisions, TierMoE achieves a \textbf{$+5.02\text{ pp}$ to $+8.68\text{ pp}$ higher hit rate} over single-request controls, outperforms our predictive lookahead heuristic by \textbf{$+6.99\text{ pp}$}, and outperforms our reactive CXL-LRU tiering heuristic by up to \textbf{$+32.27\text{ pp}$}. Furthermore, our findings show that pairwise co-activation tracking did not provide a statistically significant benefit over instantaneous batch frequency, while CXL sensitivity analysis supports a bandwidth-dominant first-order model for the tested bulk-transfer regime. TierMoE establishes an effective, low-overhead foundation for scaling MoE models on memory-tiered architectures.

% -------------------------------------------------------------
\section*{Acknowledgment}
The author thanks the Department of Electrical and Electronics Engineering at BITS Pilani, K. K. Birla Goa Campus, for computational resources and support throughout this investigation.

% -------------------------------------------------------------
\begin{thebibliography}{10}
\providecommand{\url}[1]{#1}
\csname url@samestyle\endcsname
\providecommand{\newblock}{\relax}
\providecommand{\bibinfo}[2]{#2}
\providecommand{\BIBentrySTDinterwordspacing}{\spaceskip=0pt\relax}
\providecommand{\BIBentryALTinterwordstretchfactor}{4}
\providecommand{\BIBentryALTinterwordspacing}{\spaceskip=\fontdimen2\font plus
\BIBentryALTinterwordstretchfactor\fontdimen3\font minus
  \fontdimen4\font\relax}
\providecommand{\BIBforeignlanguage}[2]{{%
\expandafter\ifx\csname l@#1\endcsname\relax
\typeout{** WARNING: IEEEtran.bst: No hyphenation pattern has been}%
\typeout{** loaded for the language `#1'. Using the pattern for}%
\typeout{** the default language instead.}%
\else
\language=\csname l@#1\endcsname
\fi
#2}}
\providecommand{\BIBdecl}{\relax}
\BIBdecl

\bibitem{jiang2024mixtral}
A.~Q. Jiang \emph{et~al.}, ``Mixtral of experts,'' \emph{arXiv preprint arXiv:2401.04088}, 2024.

\bibitem{deepseek2024deepseekv3}
DeepSeek-AI, ``DeepSeek-V3 technical report,'' \emph{arXiv preprint arXiv:2412.19437}, 2024.

\bibitem{qwen2024qwen25}
Qwen~Team, ``Qwen2.5: A party of foundation and large language models,'' \emph{arXiv preprint arXiv:2412.15115}, 2024.

\bibitem{rajbhandari2022deepspeed}
S.~Rajbhandari \emph{et~al.}, ``DeepSpeed-MoE: Advancing mixture-of-experts inference and training to unprecedented scale,'' in \emph{Proc. ICML}, 2022, pp. 18\,333--18\,346.

\bibitem{fiddler2024}
F.~Fiddler \emph{et~al.}, ``Fiddler: Serving large MoE models on a single GPU with CPU offloading,'' in \emph{Proc. EuroSys}, 2024.

\bibitem{fu2024moeinfinity}
L.~Fu \emph{et~al.}, ``MoE-Infinity: Fast MoE serving with activation-aware expert offloading,'' in \emph{Proc. USENIX OSDI}, 2024.

\bibitem{promoe2025}
X.~ProMoE \emph{et~al.}, ``ProMoE: Fast, memory-efficient MoE serving via predictive expert prefetching,'' in \emph{Proc. ACM ASPLOS}, 2025.

\bibitem{cxlmoe2023}
Y.~CXL-MoE \emph{et~al.}, ``CXL-MoE: Accelerating mixture-of-experts LLM serving via CXL memory pooling and near-data processing,'' in \emph{Proc. IEEE/ACM ISCA}, 2023.

\bibitem{hybrimoe2024}
Z.~HybriMoE \emph{et~al.}, ``HybriMoE: Hybrid memory scheduling for heterogeneous MoE inference,'' in \emph{Proc. IEEE MICRO}, 2024.

\bibitem{firmmoe2024}
K.~FIRM-MoE \emph{et~al.}, ``FIRM-MoE: Fine-grained memory management for CXL-based MoE inference,'' in \emph{Proc. IEEE HPCA}, 2024.

\bibitem{cxlmemsim2023}
CXLMemSim~Team, ``CXLMemSim: A cycle-level discrete-event simulator for CXL memory expanders,'' \emph{GitHub repository}, 2023.

\bibitem{cxlconsortium2022}
{CXL Consortium}, ``Compute Express Link Specification, Revision 3.0,'' Tech. Rep., 2022.

\bibitem{shazeer2017outrageously}
N.~Shazeer \emph{et~al.}, ``Outrageously large neural networks: The sparsely-gated mixture-of-experts layer,'' in \emph{Proc. ICLR}, 2017.

\bibitem{fedus2022switch}
W.~Fedus, B.~Zoph, and N.~Shazeer, ``Switch transformers: Scaling to trillion parameter models with simple and efficient sparsity,'' \emph{JMLR}, vol.~23, no. 120, pp. 1--39, 2022.

\bibitem{sglang2024}
L.~Zheng \emph{et~al.}, ``Efficiently scaling LLM serving with structured generation and SGLang,'' in \emph{Proc. USENIX OSDI}, 2024.

\end{thebibliography}

\end{document}
